\chapter{Trades de swap spread et trades d'asset swap}\label{ch:s2:swap-spread-and-asset-swap-trades}

D'octobre 2008 à octobre 2016, le taux de swap du dollar américain à trente ans est resté sous le rendement du Trésor de même échéance pendant 94~\% des jours. Un manuel
affirme que cela ne peut pas durer~: acheter le Trésor, payer le taux fixe dans le swap et encaisser la différence. Cela a duré parce que le trade qui devrait y mettre fin mobilise
du bilan, et que les banques ne voulaient plus prêter le leur à bon marché. Jermann a montré que, lorsque la détention d'obligations engendre des frictions, les spreads de swap
négatifs ne devraient surprendre personne. Sur le marché synthétique de ce chapitre, le \omterm{def:s2:swap-spread-and-asset-swap-trades:trade}{trade de swap spread} long rapporte 32.4 points de base par an sur son notionnel une fois
qu'un \omterm{def:s2:swap-spread-and-asset-swap-trades:bscost}{coût de bilan} a fait passer le spread en territoire négatif. Imputé au bilan qu'il utilise, il perd 17.6. La construction est \texttt{firm.swapspread}.

\section{Spreads de swap et asset swaps}

Un spread de swap (livre~2, chapitre~9) est le taux fixe d'un swap de taux d'intérêt moins le rendement de l'obligation d'État de même échéance. Un spread d'asset
swap (livre~2, chapitre~21) est la marge au-dessus du taux variable que verse une obligation une fois swappée~: acheter l'obligation et payer le taux fixe d'un swap adossé aux
coupons de l'obligation la transforme en actif à taux variable. Avant 2008, les spreads de swap étaient positifs, une prime pour le risque de crédit bancaire dans la jambe variable et pour
la liquidité des bons du Trésor. De juillet 2000 à septembre 2008, le spread à dix ans s'est établi en moyenne à 58.6 points de base, et celui à trente ans à 48.3.

\begin{definition}[Trade de swap spread]\label{def:s2:swap-spread-and-asset-swap-trades:trade}
Un \emph{trade de swap spread}\index{trade de swap spread} (swap-spread trade) associe une obligation d'État à un swap de taux d'intérêt de même échéance, de sorte que la position
gagne lorsque le spread de swap varie~: être long du spread revient à être long de l'obligation (financée en repo) et à payer le taux fixe du swap, ce qui rapporte le rendement de l'obligation au-dessus du
taux du swap, plus le taux variable au-dessus du repo, et gagne lorsque le spread s'élargit.
\end{definition}

\section{Spreads de swap négatifs}

À partir d'octobre 2008, la partie longue de la courbe est devenue négative (\cref{fig:s2:swap-spread-and-asset-swap-trades:real}). Le spread à trente ans s'est établi en moyenne à $-21.4$ points
de base jusqu'en octobre 2016, il était négatif pendant 93.9~\% de ces jours et a atteint $-62$ le 13 septembre 2016. Le spread à dix ans s'est établi en moyenne à 8.3 et était négatif
pendant 17.6~\% des jours, atteignant $-23$ le 24 février 2016. Le spread à deux ans n'est jamais passé sous zéro~: son minimum était de 2 points de base, le 20 novembre 2015.
Le modèle de Jermann valorise les swaps lorsque la détention d'obligations est coûteuse. Il reproduit les spreads négatifs sans déséquilibres de demande importants sur le marché des swaps.

\begin{omfigure}
\begin{tikzpicture}
\begin{axis}[width=11cm,height=5.4cm,xlabel={\small année},ylabel={\small swap $-$ Trésor (pb)},tick label style={font=\footnotesize},
  xmin=2000.5,xmax=2017,ymin=-70,ymax=140,grid=major,grid style={gray!25},table/col sep=comma,/pgf/number format/1000 sep={},
  legend style={font=\scriptsize,at={(0.97,0.97)},anchor=north east}]
\addplot[omDef,thick] table[x=year,y=s10]{figdata/strategies-2/12-swap-spread-and-asset-swap-trades/real.csv};
\addplot[omThm,thick] table[x=year,y=s30]{figdata/strategies-2/12-swap-spread-and-asset-swap-trades/real.csv};
\addplot[black,thin,domain=2000.5:2017] {0};
\legend{10 ans,30 ans}
\end{axis}
\end{tikzpicture}

\omcaption{Spreads de swap en dollars américains (taux de swap ICE moins le rendement du Trésor à échéance constante), moyennes mensuelles, de juillet 2000 à octobre 2016, date à laquelle la
série de swaps sur FRED s'arrête. Dérivé de FRED~; les taux de swap ne sont pas redistribués. Données~: \texttt{s2\_fetch\_swaps}.}\label{fig:s2:swap-spread-and-asset-swap-trades:real}
\end{omfigure}

\section{Le bilan comme coût}

\begin{definition}[Coût de bilan]\label{def:s2:swap-spread-and-asset-swap-trades:bscost}
Le \emph{coût de bilan}\index{coût de bilan} (balance-sheet cost) d'un trade est le rendement qu'une banque ou un dealer exige sur les fonds propres que la réglementation la contraint à détenir
contre les actifs bruts du trade, quel que soit leur risque, imputé comme un taux courant sur le notionnel~; il rend coûteux à détenir des trades à faible risque et à fort volume, comme
les positions obligataires financées en repo.
\end{definition}

Du, Tepper et Verdelhan ont mis en évidence la même force sur le marché des changes. Les écarts par rapport à la parité des taux d'intérêt couverte étaient importants, persistants et n'étaient pas expliqués par
le risque de crédit ou les coûts de transaction. Ils étaient les plus forts pour les contrats à terme qui figuraient au bilan des banques en fin de trimestre, ce qui désigne
la réglementation comme cause.

\texttt{firm.swapspread} rend le mécanisme explicite (\cref{lst:s2:swap-spread-and-asset-swap-trades:model}). Pendant trois ans, le spread à dix ans
se situe autour d'une prime de 10 points de base. À partir de l'année 3, un \omterm{def:s2:swap-spread-and-asset-swap-trades:bscost}{coût de bilan} de 45 points de base par an entre en vigueur progressivement sur un an, et le spread le suit
à la baisse~: il s'établit en moyenne à 11.5 avant, et à $-38.8$ une fois le coût intégré, et il reste négatif chaque jour. En fin de trimestre, il recule de 6 points de base supplémentaires,
de 12 en fin d'année, et un bruit à retour à la moyenne de 8 points de base court tout au long de la période (\cref{fig:s2:swap-spread-and-asset-swap-trades:synthetic}).

\begin{omfigure}
\begin{tikzpicture}
\begin{axis}[width=11cm,height=5.4cm,xlabel={\small année},ylabel={\small spread (pb)},tick label style={font=\footnotesize},
  xmin=0,xmax=10,ymin=-75,ymax=30,grid=major,grid style={gray!25},table/col sep=comma,
  legend style={font=\scriptsize,at={(0.97,0.97)},anchor=north east}]
\addplot[omThm,thin] table[x=year,y=spread]{figdata/strategies-2/12-swap-spread-and-asset-swap-trades/synthetic.csv};
\addplot[black,thick,dashed] table[x=year,y=cost]{figdata/strategies-2/12-swap-spread-and-asset-swap-trades/synthetic.csv};
\legend{spread de swap à dix ans,moins le coût de bilan}
\end{axis}
\end{tikzpicture}

\omcaption{Le spread de swap synthétique à dix ans sur dix ans~: une faible prime, puis un \omterm{def:s2:swap-spread-and-asset-swap-trades:bscost}{coût de bilan} qui entre en vigueur à partir de l'année 3, avec des creux de fin de trimestre
et du bruit. Données~: \texttt{s2\_swapspread.market}.}\label{fig:s2:swap-spread-and-asset-swap-trades:synthetic}
\end{omfigure}

\section{Trades et leur carry}

La position sur spread long porte sur une unité de notionnel~: être long du Trésor à dix ans en repo et payer le taux fixe d'un swap à dix ans. Son DV01 est de 8.5 points de base
de notionnel par point de base, et le taux variable reçu dépasse le repo de 5 points de base par an. La charge de bilan suppose des fonds propres de 5~\% du
notionnel du Trésor, rémunérés à 10~\%, soit 50 points de base par an.

\begin{center}
\small
\begin{tabular}{@{}p{4.8cm}>{\centering\arraybackslash}p{2.5cm}>{\centering\arraybackslash}p{2.5cm}>{\centering\arraybackslash}p{2.5cm}>{\centering\arraybackslash}p{2.5cm}@{}}
\toprule
pb de notionnel par an & avant, sans charge & avant, avec charge & après, sans charge & après, avec charge\\
\midrule
carry (rendement de l'obligation moins swap) & $-11.5$ & $-11.5$ & 38.8 & 38.8\\
réévaluations & $-17.8$ & $-17.8$ & $-11.4$ & $-11.4$\\
variable au-dessus du repo & 5.0 & 5.0 & 5.0 & 5.0\\
charge de bilan & 0.0 & $-50.0$ & 0.0 & $-50.0$\\
total & $-24.3$ & $-74.3$ & 32.4 & $-17.6$\\
ratio de Sharpe & $-0.15$ & $-0.44$ & 0.13 & $-0.07$\\
\bottomrule
\end{tabular}
\end{center}

Une fois le coût intégré, le trade présente un carry positif pour qui ne paie pas le bilan~: 38.8 points de base par an issus du spread négatif.
Pour une banque qui doit détenir des fonds propres contre le Trésor, la charge absorbe tout et davantage. Le spread se fixe là où le détenteur marginal est indifférent,
et il y reste parce que le détenteur marginal est une banque. Le ratio de Sharpe est faible même sans la charge~: le bruit du spread, valorisé à un DV01 de
8.5, éclipse le carry d'une année sur l'autre.

Les creux de fin de trimestre relèvent d'un autre trade. Acheter le spread au dernier cours de clôture d'un trimestre et le vendre cinq jours plus tard a rapporté en moyenne 103.3 points de
notionnel aux cinq fins d'année qui ont suivi l'intégration du coût (le pire, 90.6), et 43.9 aux dix-huit autres fins de trimestre (le pire
$-4.3$). C'est une rémunération pour la fourniture de bilan les jours où les régulateurs le mesurent. Elle est intégrée au modèle, et sur le marché réel elle
n'est que de la taille des creux.

\section{Fiches de stratégie}

\begin{strategyfile}{Achat de spread de swap}\label{strat:s2:swap-spread-and-asset-swap-trades:long}
\sfield{Qui paie, et pourquoi} Les banques qui ne peuvent pas détenir d'obligations à bon compte~; les receveurs de swaps (fonds de pension, assureurs) qui font baisser les taux de swap.
\sfield{Instruments et plateformes} Bons du Trésor américain financés en repo~; swaps de taux compensés.
\sfield{Signal} Le spread rapporté au \omterm{def:s2:swap-spread-and-asset-swap-trades:bscost}{coût de bilan} propre au détenteur.
\sfield{Dimensionnement et exécution} Appariés en DV01~; dimensionnés selon la volatilité du spread.
\sfield{Coûts} Le \omterm{def:s2:swap-spread-and-asset-swap-trades:bscost}{coût de bilan}, le repo et la marge de compensation des swaps.
\sfield{Comment elle s'éteint} Des spreads qui deviennent encore plus négatifs lorsque le bilan se renchérit.
\sfield{Horizon, capacité, infrastructure} Des années~; lignes de repo et compensation des swaps.
\sfield{Backtest honnête} Imputer le \omterm{def:s2:swap-spread-and-asset-swap-trades:bscost}{coût de bilan} au coût réel du détenteur.
\sfield{Sources} Jermann (2020)~; spreads FRED~: le spread à trente ans négatif 93.9~\% des jours d'octobre 2008 à octobre 2016~; ce chapitre~: 32.4 points de
base par an sans imputation, $-17.6$ avec imputation.
\end{strategyfile}

\begin{strategyfile}{Carry d'asset swap}\label{strat:s2:swap-spread-and-asset-swap-trades:assetswap}
\sfield{Qui paie, et pourquoi} Les émetteurs et investisseurs qui laissent les obligations se traiter à bon marché face aux swaps.
\sfield{Instruments et plateformes} Obligations d'État ou d'agences, avec un swap adossé.
\sfield{Signal} Le spread d'asset swap rapporté au coût de financement.
\sfield{Dimensionnement et exécution} Obligation et swap appariés flux à flux.
\sfield{Coûts} Financement et \omterm{def:s2:swap-spread-and-asset-swap-trades:bscost}{coût de bilan}.
\sfield{Comment elle s'éteint} Le coût de financement de l'obligation dépasse son spread d'asset swap.
\sfield{Horizon, capacité, infrastructure} Jusqu'à l'échéance ou sur plusieurs années.
\sfield{Backtest honnête} Financement au taux du détenteur, et non au taux de l'indice variable.
\sfield{Sources} Aucun chiffre de performance vérifié.
\end{strategyfile}

\begin{strategyfile}{Trade de courbe de spread de swap}\label{strat:s2:swap-spread-and-asset-swap-trades:curve}
\sfield{Qui paie, et pourquoi} Une demande concentrée sur une échéance~: les receveurs de long terme à trente ans, et la couverture bancaire aux échéances plus courtes.
\sfield{Instruments et plateformes} Positions sur spread à deux échéances.
\sfield{Signal} La courbe des spreads (spread à trente ans moins spread à dix ans) comparée à son historique.
\sfield{Dimensionnement et exécution} Appariés en DV01 à chaque échéance.
\sfield{Coûts} Quatre jambes.
\sfield{Comment elle s'éteint} Un changement structurel d'un côté (réglementation, demande des fonds de pension).
\sfield{Horizon, capacité, infrastructure} Quelques mois à plusieurs années.
\sfield{Backtest honnête} La rupture de 2008~: les spreads à trente ans et à dix ans se sont écartés pendant des années.
\sfield{Sources} Spreads FRED (ce chapitre)~; aucun chiffre de performance vérifié.
\end{strategyfile}

\begin{strategyfile}{Trade de bilan de fin d'année}\label{strat:s2:swap-spread-and-asset-swap-trades:yearend}
\sfield{Qui paie, et pourquoi} Les banques qui réduisent leur bilan aux dates de mesure de fin de trimestre et de fin d'année.
\sfield{Instruments et plateformes} Positions sur spread, repo, swaps de change à cheval sur le passage de fin de période.
\sfield{Signal} Le calendrier~: positions prises au creux, soldées après la date.
\sfield{Dimensionnement et exécution} Uniquement avec un bilan qui n'est pas lui-même mesuré à ces dates.
\sfield{Coûts} Financement à court terme à cheval sur le passage de fin de période.
\sfield{Comment elle s'éteint} Des règles qui calculent une moyenne sur le trimestre au lieu de mesurer à sa fin.
\sfield{Horizon, capacité, infrastructure} Quelques jours.
\sfield{Backtest honnête} Dates de mesure et règles telles qu'elles étaient.
\sfield{Sources} Du, Tepper et Verdelhan (2018) sur les écarts de CIP en fin de trimestre~; les creux introduits dans ce chapitre~: 103.3 points de base aux fins d'année.
\end{strategyfile}

\section{Tutoriel~: sous le Trésor}\label{tut:s2:swap-spread-and-asset-swap-trades:tut}

\begin{tutorial}
\textbf{Objectif.} Simuler le spread de swap avec un \omterm{def:s2:swap-spread-and-asset-swap-trades:bscost}{coût de bilan}, faire tourner le trade long sur spread avec et sans imputation du coût, et trader les fins de trimestre~;
lire les spreads réels. \textbf{État final~:} le tableau et les deux figures.
\begin{tutsteps}
\item \textbf{Spread et trade}.
\omcode{code/firm/swapspread/firm_swapspread.py}{44}{74}{Le spread avec un coût de bilan introduit progressivement, et les parties du trade long sur spread.}\label{lst:s2:swap-spread-and-asset-swap-trades:model}
\item \textbf{Fins de trimestre}.
\omcode{code/strategies-2/12-swap-spread-and-asset-swap-trades/python/s2_swapspread.py}{58}{68}{Achat du creux de fin de trimestre et revente cinq jours plus tard.}\label{lst:s2:swap-spread-and-asset-swap-trades:qe}
\item \textbf{Lancer} \texttt{s2\_fetch\_swaps.py} une fois, puis \texttt{regimes()}, \texttt{quarter\_ends()} et \texttt{fig\_swapspread.py}.
\end{tutsteps}
\textbf{Pour aller plus loin.} Faire varier le \omterm{def:s2:swap-spread-and-asset-swap-trades:bscost}{coût de bilan} selon une règle sur le ratio d'effet de levier~; ajouter un spread à trente ans avec la demande des fonds de pension~; lancer le trade
pour un investisseur dont le \omterm{def:s2:swap-spread-and-asset-swap-trades:bscost}{coût de bilan} propre est de 20 points de base.
\end{tutorial}

\section{Construction~: spreads de swap}\label{bld:s2:swap-spread-and-asset-swap-trades:swapspread}

\begin{build}
\bfield{Rôle} Un spread de swap piloté par un \omterm{def:s2:swap-spread-and-asset-swap-trades:bscost}{coût de bilan}, le carry et les réévaluations du trade long sur spread, et une imputation du \omterm{def:s2:swap-spread-and-asset-swap-trades:bscost}{coût de bilan}.
\bfield{Interface} \texttt{SwapConfig(\dots)}, \texttt{simulate\_spread(cfg)}, \texttt{long\_spread(s, cfg, start, days, charge)}.
\bfield{Règles} Le carry vaut moins le spread~; les réévaluations valent le DV01 multiplié par la variation du spread~; le \omterm{def:s2:swap-spread-and-asset-swap-trades:bscost}{coût de bilan} est le capital multiplié par le taux de rendement minimal exigé, appliqué au notionnel.
\bfield{Tests d'acceptation} \texttt{code/firm/swapspread/tests/}~: un spread plat rapporte exactement son carry~; le \omterm{def:s2:swap-spread-and-asset-swap-trades:bscost}{coût de bilan} vaut 50 points de base par an~; les
creux de fin de trimestre tombent aux bons jours.
\bfield{Pour aller plus loin} Plusieurs échéances~; des asset swaps sur des obligations individuelles~; un modèle de taux de financement.
\end{build}

\begin{omsources}
\item U.~J.~Jermann, ``Negative swap spreads and limited arbitrage'', \emph{Review of Financial Studies} 33(1), 2020.
\item W.~Du, A.~Tepper and A.~Verdelhan, ``Deviations from covered interest rate parity'', \emph{Journal of Finance} 73(3), 2018.
\item Taux de swap ICE et rendements du Trésor (H.15), via FRED.
\end{omsources}

\section{Exercices}

\begin{exercise}[$\star$]\label{exo:s2:swap-spread-and-asset-swap-trades:1}
Une banque détient un capital de 5~\% contre une position sur bons du Trésor et exige un rendement de 10~\% sur ce capital. Quel est le coût de détention de la position, en points de base par an~?
\end{exercise}

\begin{exercise}[$\star$]\label{exo:s2:swap-spread-and-asset-swap-trades:2}
Une position a un DV01 de 8.5 points de base de notionnel par point de base. Le spread s'élargit de 12 points de base. Que gagne la position longue sur spread~?
\end{exercise}

\begin{exercise}[$\star$]\label{exo:s2:swap-spread-and-asset-swap-trades:3}
Le spread vaut $-40$ points de base, la jambe variable dépasse le repo de 5 et le \omterm{def:s2:swap-spread-and-asset-swap-trades:bscost}{coût de bilan} est de 50. Quel est le carry imputé, par an~?
\end{exercise}

\begin{exercise}[$\star\star$]\label{exo:s2:swap-spread-and-asset-swap-trades:4}
Pourquoi le spread à deux ans est-il resté positif alors que celui à trente ans est devenu négatif~?
\end{exercise}

\begin{exercise}[$\star\star$]\label{exo:s2:swap-spread-and-asset-swap-trades:5}
Pourquoi le creux de fin de trimestre appuie-t-il l'explication par le \omterm{def:s2:swap-spread-and-asset-swap-trades:bscost}{coût de bilan}~?
\end{exercise}

\begin{exercise}[$\star\star$]\label{exo:s2:swap-spread-and-asset-swap-trades:6}
Qui peut exécuter avec profit le trade long sur spread, et pourquoi~?
\end{exercise}

\begin{exercise}[$\star\star\star$]\label{exo:s2:swap-spread-and-asset-swap-trades:7}
\emph{Programmation.} Lancer le trade long sur spread depuis le début de l'échantillon, sans \omterm{def:s2:swap-spread-and-asset-swap-trades:bscost}{coût de bilan}. Pourquoi perd-il, et que dit cela de l'entrée dans un
trade de carry juste avant que son facteur déclencheur change~?
\end{exercise}

\begin{exercise}[$\star\star\star$]\label{exo:s2:swap-spread-and-asset-swap-trades:8}
\emph{Trouver la faille.} «~Les swap spreads négatifs constituent un arbitrage~: acheter des bons du Trésor, payer le taux fixe et encaisser 40 points de base par an sans risque.~»
\end{exercise}

\section{Problème~: sous le Trésor}

\begin{problem}[{Problème du week-end --- swap spreads et coût de bilan}]\label{pb:s2:swap-spread-and-asset-swap-trades:1}
Le spread synthétique du chapitre, les taux de swap de FRED et les données publiques.

\medskip\noindent\textbf{Partie I --- Le spread.}
\begin{enumerate}
\item Définir le swap spread et l'asset-swap spread.
\item Pourquoi les swap spreads étaient-ils positifs avant 2008~?
\item Donner les moyennes des spreads réels avant et après octobre 2008.
\item Que Jermann a-t-il montré~?
\end{enumerate}

\medskip\noindent\textbf{Partie II --- \omterm{def:s2:swap-spread-and-asset-swap-trades:bscost}{Coût de bilan}.}
\begin{enumerate}[resume]
\item Définir le \omterm{def:s2:swap-spread-and-asset-swap-trades:bscost}{coût de bilan}.
\item Que Du, Tepper et Verdelhan ont-ils trouvé~?
\item Comment le spread synthétique réagit-il au coût~?
\item Où se trouvent les creux de fin de trimestre, et pourquoi~?
\end{enumerate}

\medskip\noindent\textbf{Partie III --- Les trades.}
\begin{enumerate}[resume]
\item Définir le trade long sur spread et ses composantes.
\item Donner son rendement avant et après, avec et sans coût.
\item Pourquoi son ratio de Sharpe est-il faible~?
\item Que le trade de fin de trimestre a-t-il rapporté~?
\end{enumerate}

\medskip\noindent\textbf{Partie IV --- Le verdict.}
\begin{enumerate}[resume]
\item Énoncer le \emph{résultat nommé}~: le rendement du trade sur swap spread avec et sans \omterm{def:s2:swap-spread-and-asset-swap-trades:bscost}{coût de bilan}.
\item Pourquoi le spread se stabilise-t-il là où il se trouve~?
\item Qu'est-ce qui ferait remonter le spread~?
\item Comment dimensionneriez-vous le trade~?
\item Comment le backtester honnêtement~?
\item Quelle fiche de stratégie convient à un fonds disposant d'un bilan bon marché~?
\item Quel lien ce chapitre entretient-il avec le chapitre~11~?
\item En une phrase~: que price un swap spread négatif~?
\end{enumerate}
\end{problem}

\section{Questions d'entretien}

\begin{interviewq}[$\star$ \iqroles{trader}]\label{iq:s2:swap-spread-and-asset-swap-trades:1}
Qu'est-ce qu'un swap spread, et que signifie-t-il lorsqu'il est négatif~?
\end{interviewq}

\begin{interviewq}[$\star\star$ \iqroles{chercheur}]\label{iq:s2:swap-spread-and-asset-swap-trades:2}
Pourquoi le taux de swap à trente ans pourrait-il rester inférieur au rendement du Trésor pendant des années~?
\end{interviewq}

\begin{interviewq}[$\star\star$ \iqroles{trader}]\label{iq:s2:swap-spread-and-asset-swap-trades:3}
Construisez un asset swap sur un bon du Trésor à dix ans. Que touchez-vous, et que payez-vous~?
\end{interviewq}

\begin{interviewq}[$\star\star$ \iqroles{risque}]\label{iq:s2:swap-spread-and-asset-swap-trades:4}
Un desk est long sur swap spreads. Quels scénarios stresseriez-vous~?
\end{interviewq}

\begin{interviewq}[$\star\star$ \iqroles{développeur}]\label{iq:s2:swap-spread-and-asset-swap-trades:5}
Comment imputeriez-vous le \omterm{def:s2:swap-spread-and-asset-swap-trades:bscost}{coût de bilan} aux desks dans le système de P\&L d'une société~?
\end{interviewq}

\begin{interviewq}[$\star\star\star$ \iqroles{chercheur}]\label{iq:s2:swap-spread-and-asset-swap-trades:6}
Montrez que le carry de la position longue sur spread est égal au rendement de l'obligation moins le taux de swap, plus le taux variable, moins le taux de repo,
et expliquez pourquoi l'écart entre taux variable et repo compte lorsque l'indice variable du swap change.
\end{interviewq}
